\uuid{bj3J}
\titre{Trigonométrie}
\auteur{bodin ; hanani ; mzari}
\organisation{exo7}
\licence{CC BY-NC-SA 4.0}
\datecreate{2020-01-20}
\niveau{L1}
\module{Analyse}
\chapitre{Fonctions circulaires et hyperboliques inverses}
\sousChapitre{Autre}
\difficulte{1}
\theme{Fonctions usuelles, trigonométrie}
\exo7qcmid{lille-1.fonctions-usuelles.005}

\contenu{
\begin{question}
\qclassification{126.00}
\qtags{motcle=trigonométrie}
Soit $A=\arcsin (\sin \frac{15\pi}{7})$, $B=\arccos (\cos \frac{21\pi}{11})$ et $C=\arctan (\tan \frac{17\pi}{13})$.  Quelles sont les assertions vraies ?

\begin{answers}
    \bad{$A=\frac{15\pi}{7}$}

    \good{$A=\frac{\pi}{7}$}

    \bad{$B=-\frac{\pi}{11}$}

    \good{$C=\frac{4\pi}{13}$}
\end{answers}
\begin{explanations}
On a : $\arcsin x \in [-\frac{\pi}{2},\frac{\pi}{2}], \forall x \in [-1,1]$,   $\arccos x \in [0,\pi], \forall x\in [-1,1]$ et $\arctan x \in ]-\frac{\pi}{2},\frac{\pi}{2}[, \forall x \in \Rr$. En utilisant la périodicité des fonctions sinus, cosinus et tangente, on obtient : $A=\frac{\pi}{7}$, $B=\frac{\pi}{11}$ et $C=\frac{ 4\pi}{13}$.
\end{explanations}
\end{question}
}
